\documentclass[../../../include/open-logic-section]{subfiles}

\begin{document}

\olfileid[fr]{sth}{replacement}{ref}
\olsection{Remplacement et réflexion}

Notre dernière tentative de justification du remplacement par
\stagesinex{} commence par un résultat profond et remarquable :
\footnote{Rappel : toute formule peut contenir des paramètres,
sauf indication explicite du contraire.}
% Ici, une notation avec barre pourrait abréger la liste
% $x_1, \ldots, x_n$.
\begin{thm}[Schéma de réflexion]\ollabel{reflectionschema}
Pour toute formule $\phi$ :
\[
\forall \alpha \exists \beta > \alpha (\forall x_1 \ldots, x_n \in
V_\beta)(\phi(x_1, \ldots, x_n) \liff \phi^{V_\beta}(x_1, \ldots, x_n))
\]
\end{thm}
\noindent
Comme dans la
\olref[sth][replacement][strength]{formularelativization},
$\phi^{V_\beta}$ résulte de la restriction de tous les
quantificateurs de $\phi$ à l'ensemble~$V_\beta$.
Intuitivement, la réflexion affirme donc ceci : si $\phi$
est vraie dans la hiérarchie entière, alors elle est vraie
dans des \emph{segments initiaux} de la hiérarchie aussi
grands qu'on le souhaite.

\citet{Montague1961} et \citet{Levy1960} ont montré que
des formulations appropriées du remplacement et de la
réflexion sont équivalentes modulo~$\Z$ : ajouter l'un ou
l'autre donne~$\ZF$. (Nous démontrons ces résultats
dans la \olref[sth][replacement][refproofs]{sec}.)
Cette équivalence permet d'espérer justifier la réflexion
et le remplacement par \stagesinex{} de la manière suivante :
si nous acceptons \stagesinex, la hiérarchie devrait être
extrêmement haute, au point que rien de ce que nous pouvons
en dire ne suffise à borner sa hauteur. Nous pouvons comprendre
cela comme l'idée que, si une phrase~$\phi$ est vraie dans la
hiérarchie entière, alors elle est vraie dans des segments
initiaux de celle-ci aussi grands qu'on le souhaite. Et c'est
précisément la réflexion.

Là encore, il semble s'agir d'une tentative réellement
prometteuse de justification intrinsèque du remplacement.
Mais il y aurait beaucoup trop à en dire ici. À vous,
maintenant, de décider si elle réussit.\footnote{Vous pourriez
toutefois poursuivre votre lecture avec
\citet[95--100]{Incurvati2020}.}

\end{document}
